% Unidad U006D. Registro dodecafásico integral y descomposición 1+5+6.

\chapter{Integral dodecaphasic register, cyclic extension, and decomposition \(1+5+6\)}
\label{chap:registro-dodecafasico-integral}

The observable period of one hundred eight positions contains twelve sectors of
nine positions. This equality of cardinalities does not suffice to describe the
temporal state: the ordered support, cyclic law, sector index, and
enriched register of transports, orientations, carries, and incidences must be
distinguished. This unit constructs the four objects and proves
the integral reconstruction of the twelve memory coordinates.

\section{Ordered support and cyclic coordinate}

For \(m\in\{1,\ldots,12\}\), define
\begin{equation}
 E_m:=\{9(m-1)+1,\ldots,9m\}.
 \label{eq:u006d-bloques-nonadicos}
\end{equation}
The ordered support and the cyclic group are
\begin{equation}
 \Gamma_{108}:=
 E_1\sqcup\cdots\sqcup E_{12},
 \qquad
 \Theta_{108}:=\mathbb Z/108\mathbb Z.
 \label{eq:u006d-gamma-theta}
\end{equation}
\(\Gamma_{108}\) preserves position, order and pertenencia to block;
\(\Theta_{108}\) preserves the law of translation. No are the same object.

For \(\theta\in\Theta_{108}\), let
\(\widetilde\theta\in\{0,\ldots,107\}\) be its representative. We define
\begin{equation}
 \beta(\theta):=
 \left[
 \left\lfloor\frac{\widetilde\theta}{9}\right\rfloor
 \right]_{12},
 \qquad
 r(\theta):=1+(\widetilde\theta\bmod9).
 \label{eq:u006d-coordenadas-bloque-residuo}
\end{equation}
The first coordinate locates the duodecaphase sector; the second publishes
the positive position \(1,\ldots,9\) within the sector.

\section{Nonsplit cyclic extension}

Let \(C_n=\mathbb Z/n\mathbb Z\). The applications
\begin{equation}
 \iota:C_{12}\longrightarrow C_{108},
 \quad \iota([b]_{12})=[9b]_{108},
 \qquad
 p:C_{108}\longrightarrow C_9,
 \quad p([\theta]_{108})=[\theta]_9
 \label{eq:u006d-mapas-extension}
\end{equation}
form the sequence
\begin{equation}
 0\longrightarrow C_{12}
 \xrightarrow{\;\iota\;}C_{108}
 \xrightarrow{\;p\;}C_9
 \longrightarrow0.
 \label{eq:u006d-extension}
\end{equation}
For the set-theoretic section
\[
 s([r]_9):=[\widetilde r]_{108},
 \qquad 0\leq\widetilde r<9,
\]
the defect of additivity is
\begin{equation}
 c(r,r'):=
 \left[
 \left\lfloor
 \frac{\widetilde r+\widetilde r'}9
 \right\rfloor
 \right]_{12},
 \qquad
 s(r)+s(r')
 =s(r+r')+\iota(c(r,r')).
 \label{eq:u006d-cociclo-extension}
\end{equation}

\begin{theorem}[Accuracy, non-division, and cohomological class]
\label{thm:u006d-extension-no-escindida}
The sucesion \eqref{eq:u006d-extension} is exacta and not it escinde. For the
action trivial,
\begin{equation}
 H^2(C_9,C_{12})
 \simeq C_{12}/9C_{12}
 \simeq C_3,
 \label{eq:u006d-cohomologia-extension}
\end{equation}
and the cocycle \eqref{eq:u006d-cociclo-extension} represents its generating
class.
\end{theorem}

\begin{proof}
The image of \(\iota\) is the subgroup of multiples of nine, which coincides
with \(\ker p\); \(p\) is surjective. This proves exactness. If the
extension were split, \(C_{108}\) would be isomorphic to
\(C_{12}\times C_9\), but
\[
 \exp(C_{108})=108,
 \qquad
 \exp(C_{12}\times C_9)=\operatorname{mcm}(12,9)=36.
\]
The cocycle identity is the Euclidean division of \(\widetilde r+\widetilde r'\) by nine.
For a cyclic group with trivial action, \(H^2(C_9,C_{12})\simeq C_{12}/9C_{12}\); the unit carry
projects to \([1]\), which generates the quotient of order three.
\end{proof}

The non-splitting typifies the return:
\[
 g_0(\theta+9)=g_0(\theta),
 \qquad
 \beta(\theta+9)=\beta(\theta)+1\pmod{12}.
\]
Nine iterations close the visible phase and shift one sector; one hundred
eight close both coordinates of the clock. The enriched state further preserves
the vertical memory, cylinder, boundary, and genealogy.

\section{Enriched temporal register}

\begin{definition}[Oriented dodecaphonic register]
\label{def:u006d-registro-dodecafasico}
For an enriched history, let \(\tau_m\) be the ordered subroute on
\(E_m\), \(\sigma_m\) its orientation, \(\kappa_m\) the carry crossing the
block boundary, and \(\Lambda_m\) its local incidence ledger. The
complete record is
\begin{equation}
 \mathcal R_{12}:=
 \bigl(
 (E_m,\tau_m,\sigma_m,\kappa_m,\Lambda_m)
 \bigr)_{m=1}^{12}.
 \label{eq:u006d-registro-r12}
\end{equation}
\end{definition}

Thus four levels are distinguished:
\begin{equation}
 \boxed{
 \mathcal R_{12}
 \neq C_{12}
 \neq\Theta_{108}
 \neq\Gamma_{108}.}
 \label{eq:u006d-cuatro-objetos}
\end{equation}
The first is an enriched register; the second, a sector index; the
third, a group of translations; the fourth, an ordered support. There exist
projections between them, but no identifications.

Let \(q_{\rm mem}\in\mathbb Z_{\geq0}\) be the unbounded number of nonadic
turns. The duodecaphase coordinate is only its residue
\[
 \beta=[q_{\rm mem}]_{12}.
\]
After one hundred and eight steps,
\[
 \beta\longmapsto\beta,
 \qquad
 q_{\rm mem}\longmapsto q_{\rm mem}+12.
\]
The calendar returns; the vertical memory is not reset.

\section{Signed event extractor}

Let \((d_t)_{t=1}^{108}\) be the digital record on
\(\Gamma_{108}\). We define
\begin{equation}
 \mathcal D_{\rm evt}:=\{2,4,6,8\},
 \qquad
 \Sigma_{12}:=(+,+,+,-,-,-,+,+,+,-,-,-).
 \label{eq:u006d-datos-eventos}
\end{equation}
For \(m=0,\ldots,11\), let
\begin{equation}
 e_m:=
 \Sigma_{12}(m+1)\,
 \#\{t\in E_{m+1}:d_t\in\mathcal D_{\rm evt}\}
 \in\mathbb Z.
 \label{eq:u006d-eventos-firmados}
\end{equation}
The application preserves the local census and the sectorial orientation as separate factors.

We pair the two semicrowns by means of
\begin{equation}
 p_i:=e_i+e_{i+6},
 \qquad
 a_i:=e_i-e_{i+6},
 \qquad 0\leq i\leq5.
 \label{eq:u006d-pares-antipodales}
\end{equation}
We define the total census, the five differences with respect to the sixth pair,
and the six antisymmetries:
\begin{equation}
 s_C:=\sum_{i=0}^{5}p_i,
 \qquad
 M_i:=p_i-p_5\quad(0\leq i\leq4),
 \qquad
 A_6:=(a_0,\ldots,a_5).
 \label{eq:u006d-coordenadas-uno-cinco-seis}
\end{equation}

\begin{definition}[Integral reader \(1+5+6\)]
\label{def:u006d-lector-integral}
The linear reader is
\begin{equation}
 \mathcal L_{12}:\mathbb Z^{12}\longrightarrow\mathbb Z^{12},
 \qquad
 \mathcal L_{12}(e)
 :=
 (s_C,M_0,\ldots,M_4,a_0,\ldots,a_5).
 \label{eq:u006d-lector-l12}
\end{equation}
The coordinate distribution is
\begin{equation}
 \underbrace{1}_{s_C}
 +\underbrace{5}_{M_0,\ldots,M_4}
 +\underbrace{6}_{a_0,\ldots,a_5}
 =12.
 \label{eq:u006d-uno-cinco-seis}
\end{equation}
\end{definition}

\begin{theorem}[Exact reconstruction and integral image]
\label{thm:u006d-inversa-l12}
The reader \(\mathcal L_{12}\) is injective. On its image,
\begin{align}
 p_5&=
 \frac{s_C-\sum_{i=0}^{4}M_i}{6},
 &
 p_i&=p_5+M_i\quad(0\leq i\leq4),
 \label{eq:u006d-inversa-p}\\
 e_i&=\frac{p_i+a_i}{2},
 &
 e_{i+6}&=\frac{p_i-a_i}{2}
 \quad(0\leq i\leq5).
 \label{eq:u006d-inversa-e}
\end{align}
An output tuple belongs to \(\operatorname{im}\mathcal L_{12}\) if and
only if
\begin{equation}
 s_C-\sum_{i=0}^{4}M_i\equiv0\pmod6
 \label{eq:u006d-condicion-divisibilidad}
\end{equation}
and, for the reconstructed \(p_i\),
\begin{equation}
 p_i\equiv a_i\pmod2
 \qquad(0\leq i\leq5).
 \label{eq:u006d-condicion-paridad}
\end{equation}
\end{theorem}

\begin{proof}
From \(M_i=p_i-p_5\) is obtained
\[
 s_C=\sum_{i=0}^{4}(p_5+M_i)+p_5
 =6p_5+\sum_{i=0}^{4}M_i,
\]
which gives the first equation of \eqref{eq:u006d-inversa-p}; the others are
immediate. Adding and subtracting
\(p_i=e_i+e_{i+6}\) and \(a_i=e_i-e_{i+6}\) produces
\eqref{eq:u006d-inversa-e}. Divisibility by six and the six
parity congruences are precisely the necessary and
sufficient conditions for all reconstructed coordinates to be integers.
\end{proof}

Division by six appears after the oriented census, not during it. The
integrality conditions for this unit are exactly
\eqref{eq:u006d-condicion-divisibilidad} and
\eqref{eq:u006d-condicion-paridad}; no Smith normal form is attributed here to a
matrix that has not previously been defined.

\section{Two dodecaphasic words of distinct origins}

The distinction must be preserved.
\begin{equation}
 U_{12}^{\rm cat}:=U_6\Vert U_6,
 \qquad
 U_{12}^{\rm sgn}:=\mathcal H_{12}(K).
 \label{eq:u006d-dos-palabras-doce}
\end{equation}
The first concatenates two catalog emissions of length six. The second
is a signed coordinate reversibly linked to \(K\) by the operator
\(\mathcal H_{12}\). Their common length does not imply the same provenance,
semantics, or transport law.

\begin{criterion}[Falsification criteria]
\label{crit:u006d-refutacion}
The construction is refuted if the sequence
\eqref{eq:u006d-extension} is not exact, if its cohomology class is
trivial, if a tuple in the image violates
\eqref{eq:u006d-condicion-divisibilidad} or
\eqref{eq:u006d-condicion-paridad}, if the inverse formulas do not recover
the twelve initial integers, or if the return of \(\beta\) is interpreted as a
reset of \(q_{\rm mem}\).
\end{criterion}
